\input{../article_base.tex}
\title{פתרון מטלה 12 --- משוואות דיפרנציאליות, 80320}
% chktex-file 9
% chktex-file 17

\DeclareMathOperator\Div{div}

\usepackage{pgfplots, tikz, pgffor}
\usetikzlibrary{math}
\pgfplotsset{compat=1.18}

\begin{document}
\maketitle
\maketitleprint[blue]{}

\question{}
תהי $\Omega \subseteq \RR^d$ פתוחה, לכל $K \subseteq \Omega$ קומפקטית ו־$f : \Omega \to \RR$ רציפה נגדיר,
\[
	\lVert f \rVert_{\infty, K}
	= \sup_{x \in K} |f(x)|
.\]
נניח ש־${(u_n)}_{n = 1}^{\infty} \subseteq C^2(\Omega)$ סדרת הרמוניות ונניח שקיימת $u : \Omega \to \RR^d$ כך שלכל $K \subseteq \Omega$ קומפקטית מתקיים,
\[
	\lim_{n \to \infty} \lVert u_n - u \rVert_{\infty, K}
	= 0
.\]

\subquestion{}
נראה ש־$u$ הרמונית.
\begin{proof}
	ניזכר תחילה כי $\RR^d$ היא $\sigma$־קומפקטית ולכן בפרט גם $\Omega$, ונסמן ${\{ K_m \}}_{m = 1}^\infty \subseteq \Tt(\RR^d)$ להיות קבוצת קומפקטיות כזו, כלומר,
	\[
		\bigcup_{m = 1}^\infty K_m
		= \Omega
	.\]
	בנוסף אנו מקבלים ש־$\lim_{n \to \infty} \lVert u_n - u \rVert_{\infty, K_m} = 0$, ולכן נוכל לבנות במעבר דרך הגדרת הגבול מטיעון אלכסון גם את,
	\[
		\lim_{n \to \infty} \lVert u_n - u \rVert_{\infty, \bigcup_{m = 1}^n K_m}
	\]
	בפרט גבול זה קיים ומתאפס.
	מצאנו אם כך ש־$u |_{\bigcup_{m = 1}^n K_m}$ היא פונקציה הרמונית כגבול במידה שווה של הרמוניות, ונסמן פונקציה זו ב־$u^n$.
	מתקיים $u \equiv \bigcup_{n = 1}^\infty u^n$ מהגדרה, וכן לכל נקודה $x \in \Omega$ קיים $N$ כך ש־$x \in \dom u^n$ לכל $n > N$ ובפרט $\Delta u(x) = \Delta u^n(x) = 0$ מוגדר וקיבלנו ש־$u$ הרמונית.
\end{proof}

\subquestion{}
תהי $C \subseteq \Omega$ קומפקטית,
נראה שלכל סדרת אינדקסים $i_1, \ldots, i_k$ מתקיים,
\[
	\lim_{n \to \infty} \lVert \partial_{i_1} \cdots \partial_{i_k} (u_n - u) \rVert_{\infty, C}
	= 0
.\]
\begin{proof}
	ממשפט החסם לנגזרות של הרמוניות קיים חסם גלובלי $D = D(d, k)$ כך שלכל הרמונית $v$ בתחום מתקיים,
	\[
		|\partial_{i_1} \cdots \partial_{i_k} v(x)|
		\le \frac{D}{r^{d + k}} \int_{B(x, r)} |v(y)|\ dy
	.\]
	לכל $r$ כך ש־$B(x, r) \subseteq C$.
	מקומפקטיות $C$ בתוך $\Omega$ קיים כיסוי סופי של כדורים $C \subseteq \bigcup_{n = 1}^l B(x_n, r_n) \subseteq \Omega$. \\
	נגדיר את $C' = \cl \bigcup_{n = 1}^l B(x_n, r_n)$, אז $C \subseteq C' \subseteq \Omega$ וכן $C'$ קומפקטית ולכן גם,
	\[
		l \int_{C'} |(u - u_n)(y)|\ dy
		< \infty
	.\]
	ונסמן ערך זה ב־$A_n$, נבחין כי זהו חסם לסכום האינטגרלים על הכדורים, בפרט $\int_C |u - u_n| \le A_n$, ומהתכנסות במידה שווה של $u_n$ נקבל שגם $A_n \to 0$.
	עתה לכל $x \in C$ נוכל לבחור את $x \in B(x_n, r_n)$ שקיים מעצם הכיסוי, ולכן $B(x, \lVert x - x_n \rVert) \subseteq B(x_n, r_n)$ ובפרט,
	\[
		|\partial_{i_1} \cdots \partial_{i_k} v(x)|
		\le \frac{D}{r^{d + k}} A_n
	.\]
	וקיבלנו התכנסות.
\end{proof}

\question{}
תהי $\Omega \subseteq \RR^d$ פתוחה וחסומה ונניח ש־$u \in C^2(\Omega) \cap C(\overline{\Omega})$ הרמונית ב־$\Omega$. \\
נראה ללא שימוש בתכונת הערך הממוצע שמתקיים,
\[
	\max_{\overline{\Omega}} u
	= \max_{\partial \Omega} u
.\]
\begin{proof}
	יהי $t > 0$ ונגדיר את,
	\[
		u_t(x)
		= u(x) + t \lVert x \rVert^2
	.\]
	אז מתקיים,
	\[
		\Delta u_t(x)
		= \Delta u(x) + t \Delta \lVert x \rVert^2 + \partial_t^2 u_t
		= 0 + t \sum_{i = 1}^d \partial_i^2 (x_1^2 + \cdots + x_d^2)
		= 2^d t
	.\]
	נבחין כי $\Omega$ חסומה ולכן $\cl \Omega$ קומפקטית ובפרט $u$ מקבלת עליה מקסימום.
	אם המקסימום שייך לשפה אז סיימנו, לכן נניח ש־$u(x) = \max_{\overline{\Omega}} u$ מקיימת $x \in \Omega$.
	נבחין כי $|H u|$ חסום ובפרט עבור בחירת $t$ מספיק גדול קיבלנו שהוא יהיה חיובי לחלוטין בכל נקודה, בפרט לא תהיה נקודת מקסימום מקומי.
	אם כך נובע ש־$|H u_t(x)|$ מתאפסת לאיזשהו $t > 0$, כל זה כמובן לחלוטין בלתי־אפשרי בהינתן ש־$\Delta u_t(x) = 2^d t > 0$ תמיד.
\end{proof}

\question{}
תהי $\Omega \subseteq \RR^d$ פתוחה וחסומה.
לכל $g \in C(\partial \Omega)$ נגדיר את $T g : \overline{\Omega} \to \RR$ להיות הפונקציה היחידה ב־$C^2(\Omega) \cap C(\overline{\Omega})$ שפותרת את הבעיה,
\[
	\begin{cases}
		\Delta u|_{\Omega} = 0 \\
		u |_{\partial \Omega} = g |_{\partial \Omega}
	\end{cases}
.\]
נראה ש־$T : C(\partial \Omega) \to C^2(\Omega) \cap C(\overline{\Omega})$ הוא אופרטור לינארי חסום וכן ש־$\lVert T \rVert_{\operatorname{op}} \le 1$.
\begin{proof}
	תהי $g \in C(\partial \Omega)$ ונבחין כי אם $\alpha \in \RR$ סקלר אז,
	\[
		\Delta (\alpha u)
		= \alpha \Delta u
	.\]
	מלינאריות הנגזרת.
	בפרט על $\Omega$ נקבל $\Delta(\alpha u) = \alpha \Delta u = \alpha \cdot 0 = 0$ וכן על $\partial u$,
	\[
		\alpha u = \alpha g
	.\]
	ישירות מ־$u = g$, נסיק ש־$\alpha u$ הפתרון היחיד של הבעיה עבור $\alpha g$ ולכן $T(\alpha g) = \alpha T g$.

	אם $h \in C(\partial \Omega)$ בנוסף, אז שוב מלינאריות הנגזרת נקבל גם,
	\[
		T(g + h)
		= Tg + Th
	.\]
	וקיבלנו ש־$T$ לינארי.

	נעבור לבדוק שהוא חסום,
	\[
		\lVert T g \rVert_{\infty}
		= \sup\{ |g(x)| \mid x \in \overline{\Omega} \}
	.\]

	TODO
\end{proof}

\end{document}
